\documentclass[11pt]{article}
\usepackage{fontspec}
\usepackage[french]{babel}
\usepackage{amsmath}
\usepackage{amssymb}
\usepackage{amsthm}
\usepackage{mathrsfs}
\usepackage{tikz-cd}
\usepackage{enumitem}

\newcommand{\printedpage}[1]{\marginpar{\footnotesize p.~#1}}
\newcommand{\msp}[1]{\marginpar{\footnotesize [#1]}}
\newcommand{\uHom}{\underline{\mathrm{Hom}}}
\newcommand{\Ind}{\mathrm{Ind}}
\newcommand{\Ob}{\operatorname{Ob}}
\newcommand{\Fl}{\operatorname{Fl}}
\newcommand{\card}{\operatorname{card}}
\newcommand{\cardred}{\mathop{\textup{card\,réd}}\nolimits}
\newcommand{\cM}{\mathcal{M}}
\newcommand{\rbox}[1]{\tikz[baseline=(X.base)]\node[draw,rounded corners=4pt,inner sep=2pt](X){\textup{#1}};}

\begin{document}

\printedpage{79}
% head: A. Grothendieck : Dérivateurs, Ch. XVIII, Édition des Universités de Montpellier II et Paris VII

Le résultat principal dans la présente section est le suivant :

\textbf{Théorème 4.9.3.} \textit{Sous les conditions générales de 4.9.1 pour $C$, $\pi_0$, $J$ (notamment} $\cardred J \leqslant \pi_0$\textit{), le foncteur pleinement fidèle (4.9.3) est une équivalence de catégories dans chacun des deux cas suivants :}
\begin{itemize}[label={}]
{\itshape
\item[a)] $C$ est stable par $\varinjlim$ filtrantes dénombrables.
\item[b)] Il existe un cardinal $\pi_0 < \pi$, et un cardinal $\pi'_0 > \pi_0$ \textup{(}\textbf{\upshape N.B.} On peut supposer que $\pi'_0$ est le successeur de $\pi_0$, donc $\pi'_0 \leqslant \pi$), tels que $C$ soit stable par limites inductives grandes devant $\pi_0$, et de cardinal $\pi'_0$.
}
\end{itemize}
(\footnote[94]{Cas commun : il existe un ordinal limite $\alpha$, tel que 1) $\card \alpha \leqslant \pi_0$. 2) Les limites inductives \texttt{[de]} type $I_\alpha = \{\beta \mid \beta < \alpha\}$ sont représentables dans $C$.}).

\textbf{Remarque 4.9.4.}
\begin{itemize}[label={}]
\item[a)] Ce théorème, cas a), contient le théorème 4.4.13, partie b) (p. 95-96), où on suppose que $C$ est stable par limites inductives de cardinal $\leqslant \pi_0$, l'hypothèse a) de 4.9.3 étant considérablement plus faible.
\item[b)] \msp{page 220}Je présume qu'en dehors d'une hypothèse restrictive du type a) ou b) ci-dessus, le foncteur pleinement fidèle (4.9.3) n'est pas toujours essentiellement surjectif, même dans le cas particulier où $J$ est une catégorie finie, p.ex. $J = B_G$, $G = \{\pm 1\}$. Pour un contre-exemple dans le cas très voisin où on prend $\cM = \Ind(C)$ (au lieu de $\Ind_{\pi_0}(C)$), et où on considère le foncteur (pleinement fidèle si $J$ est une catégorie équivalent à une catégorie \textit{finie}) % [sic] équivalent -> équivalente, agreement with « une catégorie »
\[
\Ind(\uHom(J,C)) \longrightarrow \uHom(J,\Ind(C)),
\]
je renvoie à l'exercice SGA 4 I 8.8.8.
\end{itemize}
L'application principale du théorème 4.9.3 consiste en le

\textbf{Corollaire 4.9.5.} \textit{Soient $\cM$ une catégorie} accessible \textit{(4.8.7), $J$ une catégorie équivalente à une petite catégorie. Alors la catégorie $\cM(J^\circ) = \uHom(J,\cM)$ est accessible. Plus précisément, soient $\pi_0$ un cardinal} utile \textit{pour $\cM$ (i.e. tel que $\Ind_{\pi_0}(\cM_{\pi_0}) \xrightarrow{\sim} \cM$, i.e. un cardinal tel que $\cM$ soit $\pi_0$-accessible), soit $\pi$ un cardinal infini tel que $\pi^{\pi_0} = \pi \geqslant \cardred J$} (\footnote[95]{On a donc $\pi = \pi^{\pi_0} \geqslant 2^{\pi_0} > \pi_0$.}), \textit{i.e. $\pi = c^{\pi_0}$, où $c$ est n'importe quel cardinal $c \geqslant \sup(\pi_0, \cardred J)$, alors $\uHom(J,\cM)$ est $\pi$-accessible, plus précisément on a une équivalence de catégories} % ?? uncertain: font of « alors » -- italic or upright
\[
\uHom(J,\cM) \xleftarrow{\ \sim\ } \Ind_\pi(\uHom(J,\cM_\pi)) \ . \tag*{(4.9.5.1)}
\]

\msp{page 221}On sait en effet, avec les hypothèses faites sur $\pi_0$, $\pi$, que $\pi$ est utile pour $\cM$, i.e. que
\[
\cM \xleftarrow{\ \sim\ } \Ind_\pi(\cM_\pi) \ ,
\]
d'autre part, $\cM_\pi$ est stable par $\varinjlim$ grandes devant $\pi_0$ et de cardinal $\leqslant \pi$. On peut donc appliquer 4.9.3 b), en prenant $C = \cM_\pi$, $\pi$ et $\pi_0$ comme ici (donc $\pi_0 < 2^{\pi_0} \leqslant \pi$), et $\pi'_0 = \pi$.

\printedpage{80}
% head: A. Grothendieck : Dérivateurs, Ch. XVIII, Édition des Universités de Montpellier II et Paris VII

\textbf{4.9.6} (\footnote[96]{Résultats préliminaires (à la démonstration de 4.9.3).}). Le reste de la présente section 4.9 est consacré à la démonstration de 4.9.3, i.e. à la vérification, pour $J \longrightarrow \cM$, des conditions a), b) de 4.9.2. Dans toute la suite, on se fixe donc une
\[
\varphi \ : \ J \longrightarrow \cM \ = \ \Ind_\pi(C) \ . \tag*{(4.9.6.1)}
\]
On peut supposer $J$ réduite, donc petite et $\card J \leqslant \pi_0$. On va considérer le foncteur déduit du foncteur (4.9.2.1)
\[
\uHom(J,C)/\varphi \longrightarrow \prod_{j \in \Ob J} C/\varphi(j) \ , \tag*{(4.9.6.2)}
\]
dont les foncteurs en question se déduisent en composant avec les foncteurs projection
\[
\prod C/\varphi(j) \longrightarrow C/\varphi(j_0) \ . \tag*{(4.9.6.3)}
\]
Comme les catégories facteurs $C/\varphi(j)$ sont filtrantes (et même grandes devant $\pi$), leur produit est également filtrant, et même grand devant $\pi$ (immédiat), et de même
\msp{page 222}
pour les produits partiels. Il s'ensuit aussitôt que les projections (4.9.6.3) sont cofinales, en vertu du

\textbf{Lemme 4.9.6.1.} \textit{Soient $I$, $I'$ deux catégories avec $I'$ \textup{0}-connexe (p.ex. filtrante), alors le foncteur projection $pr_1 : I \times I' \longrightarrow I$ est cofinal.}

En effet, pour $i' \in \Ob I'$ \texttt{[plutôt} $i \in \Ob I$\texttt{]}, on a
\[
i\backslash(I \times I') \ \simeq \ (i\backslash I) \times I' \ ,
\]
or $i\backslash I$ et $I'$ étant 0-connexe, il en est de même de leur produit.

\textbf{Lemme 4.9.6.2} (\footnote[97]{Attention, il n'est pas vrai qu'un produit infini de catégories 0-connexes est 0-connexe !}). \textit{Soit $(I_\alpha)_{\alpha \in A}$ une famille de catégories} filtrantes \textit{(resp. grandes devant $\pi$). Alors $I_A \stackrel{\text{déf}}{=} \prod_{\alpha \in A} I_\alpha$ est filtrante (resp. grande devant $\pi$), et pour toute partie $B \subset A$, la projection canonique $I_A \longrightarrow I_B$ est cofinal.} % [sic] cofinal -> cofinale, agreement with « la projection »

(\footnote[98]{Cf. démonstration élégante ci-dessous, 4.9.6.6.}). La stabilité de la propriété “filtrante” ou “grande devant $\pi$” par AQT \texttt{[âne qui trotte]}. On a, si $B' = A\backslash B$, $I_A \simeq I_B \times I_{B'} \longrightarrow I_B$, et on applique le lemme 4.9.6.1 pour conclure que ce foncteur est cofinal, $I_{B'}$ étant filtrante, donc 0-connexe.

\msp{page 223}Le foncteur (4.9.6.3) étant cofinal, pour établir que (4.9.2.1) l'est, il suffit de prouver qu'il en est ainsi du foncteur (4.9.6.2), ce que nous ferons. Notons de plus que ledit foncteur est \textit{fidèle}, car il s'insère dans le carré commutatif
\[
\begin{tikzcd}[column sep=large, row sep=large]
\uHom(J,C)/\varphi \arrow[r] \arrow[d] & \prod C/\varphi(j) \arrow[d] \\
\uHom(J,C) \arrow[r] & C^{\Ob J}\ ,
\end{tikzcd}
\]
\printedpage{81}
% head: A. Grothendieck : Dérivateurs, Ch. XVIII, Édition des Universités de Montpellier II et Paris VII
où les flèches verticales et la deuxième flèche horizontale sont fidèles. En définitive, pour établir le théorème 4.9.3, nous allons étudier le foncteur \textit{fidèle} (4.9.6.2)
\[
\begin{aligned}
&f \ : \ \uHom(J,C)/\Phi \ \longrightarrow \ \textstyle\prod_{j \in J_0} C/\varphi(j) \\ % [sic] \Phi -> \varphi, the functor is written \varphi everywhere else
&\text{(où on a posé } J_0 = \Ob J,\ J_1 = \Fl J) \ ,
\end{aligned}
\tag*{$(*)$}
\]
et nous allons prouver :
\begin{itemize}[label={}]
\item[\textcircled{a}] Ce foncteur fidèle $(*)$ est \textit{cofinal.}
\item[\textcircled{b}] La catégorie source $\uHom(J,C)/\varphi$ est grande devant $\pi_0$.
\end{itemize}
Nous allons pour ceci dégager un lemme auxiliaire sur la cofinalité des foncteurs \textit{fidèles} :

\msp{page 224}\textbf{Lemme 4.9.6.3} (\footnote[99]{Ce lemme fait partie du sorite général sur la cofinalité et les catégories filtrantes.}). \textit{Soit $f : I \longrightarrow I'$ un foncteur, et considérons les conditions :}
{\itshape
\begin{itemize}[label={}]
\item[a)] $I$ filtrante, $f$ cofinal, i.e. (en présence de $I$ filtrante) $f$ satisfait les conditions familières \rbox{F 1}, \rbox{F 2}.
\item[b)] $I'$ est filtrante, et les $i'\backslash I$ (pour $i' \in \Ob I'$) sont filtrantes.
\item[b$'$)] $I'$ \textup{0}-connexe, satisfait PF (2) \textup{(}\footnote[100]{p.ex. PF 2 est satisfaite si $I'$ est \textit{ordonnée}, sans avoir à supposer d'avance $I'$ filtrante, seulement 0-connexe.}\textup{)} (égalisation des double-flèches).
\item[b$''$)] $I'$ est filtrante, les $i'\backslash I$ pseudo-filtrantes, i.e. satisfont à la condition PF $1'$.
\item[c)] $I'$ filtrante, et $f$ satisfait à F 1, ainsi que la condition suivante.
\begin{itemize}[label={}]
\item[\rbox{F 3}] Pour tout couple $(i_0,i_1)$ d'objets de $I$ et toute flèche $f(i_0) \xrightarrow{u'} f(i_1)$, il existe un objet $i_2$ de $I$, et une flèche $v' : f(i_1) \longrightarrow f(i_2)$, tels que $v'$ et $v'u'$ soient relevables en des flèches $i_1 \xrightarrow{v} i_2$ et $i_0 \xrightarrow{w} i_2$ (i.e. il existe $i_2$ et $v : i_1 \longrightarrow i_2$, $w : i_0 \longrightarrow i_2$ tels qu'on ait $f(v)u' = f(w)$).
\end{itemize}
\end{itemize}
\[
\begin{tikzcd}
i_0 \arrow[rr, dotted, bend left=20, "w"] & i_1 \arrow[r, "v"] & i_2 % ?? uncertain: line style of the arrow v -- solid or densely dotted
\end{tikzcd}
\]
\[
\begin{tikzcd}
f(i_0) \arrow[r, "u'"] \arrow[rr, dotted, bend right=20, "f(w)"'] & f(i_1) \arrow[r, dotted, "f(v) = v'"] & f(i_2)
\end{tikzcd}
\]
\begin{itemize}[label={}]
\item[1)] On a les implications
\[
a \iff b \iff b' \implies b'' \iff c \ ,
\]
et si $f$ est \textup{fidèle} (p.ex. $I$ ordonnée), les cinq conditions a), b), b$'$), b$''$), c) sont équivalentes.
\end{itemize}
}

\printedpage{82}
% head: A. Grothendieck : Dérivateurs, Ch. XVIII, Édition des Universités de Montpellier II et Paris VII
{\itshape
\begin{itemize}[label={}]
\item[2)] \msp{page 225}Soit $\pi$ un cardinal infini, et considérons les conditions :
\begin{itemize}[label={}]
\item[$a_\pi$)] $I$ grande devant $\pi$, $f$ cofinal, i.e. F 1, F 2.
\item[$b_\pi$)] $I'$ grande devant $\pi$, les $i'\backslash I$ ($i' \in \Ob I'$) aussi.
\item[$b_\pi{}''$)] $I'$ grande devant $\pi$, les $i'\backslash I$ ($i' \in \Ob I'$) sont $\pi$-pseudo-filtrantes, i.e. satisfont PS $1_\pi{}'$. % [sic] PS 1_π' -> PF 1_π', the conditions are elsewhere called PF
\item[$c_\pi$)] $I'$ grande devant $\pi$, et $f$ satisfait F 1 et
\begin{itemize}[label={}]
\item[\rbox{F $3_\pi$}] \textup{(}\footnote[101]{\textbf{N.B.} La condition F $3_\pi$ est tautologiquement satisfaite si $f$ est \textit{pleinement fidèle.}}\textup{)}. Pour toute famille $(i_\alpha)_{\alpha \in A}$ d'objets de $I$, avec \textup{card} $A \leqslant \pi$, et $i \in \Ob I$, et une famille de flèches $u'_\alpha : f(i_\alpha) \longrightarrow f(i)$, existe $j \in \Ob I$ et une flèche $v' : f(i) \longrightarrow f(j)$, telle que $v'$ \textup{et} les $v'u'_\alpha$ proviennent de flèches de $I$, i.e. il existe des flèches $v : i \longrightarrow j$, $w_\alpha : i_\alpha \longrightarrow j$ telles que $f(v)u'_\alpha = f(w_\alpha)$ $\forall\ \alpha \in A$.
\end{itemize}
\end{itemize}
\[
\begin{tikzcd}[row sep=huge]
i_\alpha \arrow[d, dash] \arrow[rr, dotted, bend left=20, "w_\alpha"] & i \arrow[r, dotted, "v"] \arrow[d, dash] & j \arrow[d, dash] \\
f(i_\alpha) \arrow[r, "u'_\alpha"] & f(i) \arrow[r, dotted, "v' = f(v)"] & f(j)
\end{tikzcd}
\]
Ceci posé, on a
\[
a_\pi \iff b_\pi \implies c_\pi \implies b_\pi{}'' \ ,
\]
et si $f$ est fidèle, ou $I$ ordonné, les quatre conditions $a_\pi$, $b_\pi$, $b_\pi{}''$, $c_\pi$ sont équivalentes.
\end{itemize}
}
(\footnote[102]{Démonstration de la proposition donnée dans 4.9.6.7, page 227.}).

\textbf{4.9.6.4. N.B.} Dans le cas qui nous intéresse, où $f$ est le foncteur $(*)$ (p. 223), on sait que $I'$ est grande devant $\pi$. On voudrait établir que $I$ l'est et que $f$ est cofinal, i.e. la condition $\mathrm{a}_\pi$. On le prouvera en prenant $\mathrm{b}_\pi$ (qui implique $\mathrm{a}_\pi$ puisque $f$ est fidèle). Il reste donc à
\msp{page 226}
prouver, sur le foncteur $f$ dans $(*)$, p. 223, les deux conditions suivantes :
\begin{itemize}[label={}]
\item[F 1)] Pour tout système de flèches dans $\cM$
\[
X(j) \xrightarrow{p(j)} \varphi(j) \ , \qquad j \in J_0 = \Ob J, \text{ les } X(j) \text{ dans } \Ob C \ ,
\]
il existe un foncteur $F \ : \ J \longrightarrow C$ et $F \xrightarrow{q} \varphi$, enfin un système de flèches $X(j) \xrightarrow{u(j)} F(j)$ dans les $C/\varphi(j)$, i.e. rendant commutatif les diagrammes % [sic] commutatif -> commutatifs, agreement with « les diagrammes »
\[
\begin{tikzcd}[column sep=small]
X(j) \arrow[rr, "u(j)"] \arrow[dr, "p(j)"'] & & F(j) \arrow[dl, "q(j)"] \\
& \varphi(j)\ . &
\end{tikzcd}
\]
\item[F $3_\pi$)] Pour tout système $(F_\alpha)_{\alpha \in A}$ ($\card A \leqslant \pi$) d'éléments dans $\uHom(C,J)$, et des flèches % [sic] \uHom(C,J) -> \uHom(J,C), the F_α are functors J -> C
\[
F_\alpha \xrightarrow{p_\alpha} \varphi \qquad \text{dans } \uHom(J,\cM),
\]
et pour $G$ dans $\uHom(J,C)$ et $G \xrightarrow{q} \varphi$ dans $\uHom(J,\cM)$, étant donné un système de flèches $u_\alpha(j)$,
\printedpage{83}
% head: A. Grothendieck : Dérivateurs, Ch. XVIII, Édition des Universités de Montpellier II et Paris VII
\[
\begin{tikzcd}[column sep=small]
F_\alpha(j) \arrow[rr, "u_\alpha(j)", "\text{comm.}"'] \arrow[dr, "p_\alpha(j)"'] & & G(j) \arrow[dl, "q(j)"] \\
& \varphi(j)\ , &
\end{tikzcd}
\qquad u_\alpha(j) \text{ flèche dans } C/\varphi(j) \ ,
\]
il existe $H \xrightarrow{q} \varphi$ dans $\uHom(J,C)$ et une flèche $G \xrightarrow{v} H$ dans $\uHom(J,C)/\varphi$, telle que pour tout $\alpha \in A$, le système de flèches composés
\msp{page 227}
\[
\begin{tikzcd}[row sep=huge]
F_\alpha(j) \arrow[r, "u_\alpha(j)", "\text{comm.}"'] \arrow[dr, dotted, "p_\alpha(j)"'] & G(j) \arrow[r, "v(j)", "\text{comm.}"'] \arrow[d, dotted, "q(j)"] & H(j) \arrow[dl, dotted, "q(j)"] \\
& \varphi(j) &
\end{tikzcd}
\]
(pour $j \in J_0 \in \Ob J$) provienne d'un \textit{foncteur} $w_\alpha : F_\alpha \longrightarrow H$. % [sic] J_0 \in \Ob J -> J_0 = \Ob J, J_0 was defined as Ob J
\end{itemize}
Pour prouver 4.9.3, nous sommes donc ramené à prouver les deux propositions F 1), F $3_\pi$) précédentes, sous l'une des deux conditions a) ou b) de 4.9.3.

\textbf{Corollaire 4.9.6.5} (de 4.9.6.3). \textit{Soit $(f_\alpha : I_\alpha \longrightarrow I'_\alpha)$ une famille de foncteurs cofinaux, les catégories source étant filtrantes. Alors le foncteur produit $f = \prod f_\alpha : I = \prod_\alpha I_\alpha \longrightarrow I' = \prod_\alpha I'_\alpha$ est cofinal.}

En effet, il faut vérifier que les catégories $i'\backslash I$, pour $i' \in \Ob I'$, sont 0-connexes. Or si $i' = (i'_\alpha)_{\alpha \in A}$, on a
\[
i'\backslash I \ \simeq \ \prod_\alpha (i'_\alpha\backslash I_\alpha) \ ,
\]
et par 4.9.6.3, 1°) a $\Longrightarrow$ c, les $i'_\alpha\backslash I_\alpha$ sont filtrantes, donc leur produit est filtrant
\msp{page 228}
(cf. 4.4.6.2), donc 0-connexe, q.e.d. % [sic] 4.4.6.2 -> 4.9.6.2, the lemma on products of filtrantes is 4.9.6.2

\textbf{Corollaire 4.4.6.6.} \textit{Soit $\pi$ un cardinal infini. Tout produit de catégories grandes devant $\pi$ est grand devant $\pi$.} % [sic] 4.4.6.6 -> 4.9.6.6, numbering of the section and footnote 98

Soit $I = \prod_{\alpha \in A} I_\alpha$, soit $B$ un ensemble de cardinal $\leqslant \pi$, il faut prouver que $I \xrightarrow{\mathrm{diag}_I} I^B$ est cofinal. Or ce foncteur s'identifie au produit des foncteurs similaires $I_\alpha \longrightarrow I_\alpha^B$, lesquels sont cofinaux, les catégories source étant de plus filtrantes. Donc on peut conclure par le corollaire 4.9.6.5.

\textbf{4.9.6.7.} \textsc{Démonstration de} 4.9.6.3. (Partie 1°.)

a $\Longrightarrow$ b. On sait déjà que a) implique que $I'$ est filtrante (SGA 4 I 8.1.3 b). Reste à prouver que pour tout $i' \in \Ob I'$, $i'\backslash I$ est filtrante (\footnote[103]{Court-circuit : on sait déjà que $i'\backslash I$ est 0-connexe, d'autre part elle est \textit{localement} filtrante, car coétale sur $I$ qui est filtrante, donc elle est filtrante.}). Deux conditions PF $1'$ et PF 2 à vérifier (en plus du fait que $i'\backslash I \neq \emptyset$, qui est connu par cofinalité de $f$, impliquant que $i'\backslash I$ est 0-connexe).
\begin{itemize}[label={}]
\item[1)] Soient $(i,u')$, $(j,v')$ deux éléments de $i'\backslash I$, il faut voir qu'ils sont majorés par un troisième élément. Or $I$ étant filtrante, $i$, $j$ sont majorés par $k$. On aurait fini donc si on savait que $f(u_1)u' = f(v_1)v'$, soit $w'$, car alors $(k,w')$ ferait l'affaire.
\end{itemize}

\end{document}
